\chapter{Sistem ve Yöntem}
\label{ch:sistem-yontem}

\section{Donanım ve Araçlar}

Tüm cihaz-üzeri ölçümler \textbf{NVIDIA Jetson AGX Orin}~\cite{nvidia2024orin}
üzerinde alınmıştır. Orin'in CPU ve GPU'su tek bir paylaşımlı bellek
baytı ve bant genişliği üzerinden çalışır -- bu, \cref{sec:reprod-tehlike}'de
tekrar karşılaşılacak önemli bir kısıttır. Çıkarım motorları
\textbf{TensorRT}~\cite{nvidia2024tensorrt} ile derlenmiş, fp16
hassasiyetinde ve CUDA Graph tekrar oynatmasıyla çalıştırılmıştır.

\section{Kare Başına Hesaplama ve Dört Motor Kararı}
\label{sec:dort-motor}

\Cref{sec:proje-eslesme}'de gösterildiği gibi kare başına dört modül
çalışır: görüntü kodlayıcı, bellek dikkati, bellek kodlayıcı, SAM
başlığı. Bunların \emph{neden} tek bir birleşik TensorRT motoru değil
de dört ayrı motor olarak bırakıldığı gelişigüzel bir karar değildir --
modüller, aralarına \textbf{Python'un girdiği} noktalardan ayrılmıştır:

\begin{itemize}
  \item \textbf{Görüntü kodlayıcı → bellek dikkati} arasında bellek
        bankası derlenir. Bu banka kare indeksiyle anahtarlanmış bir
        Python \texttt{dict}'tir; hangi geçmiş karelerin okunacağı her
        karede \texttt{.get()} çağrısı ve en yakın adayları sıralayan bir
        fonksiyonla (\texttt{select\_closest\_cond\_frames}) belirlenir.
        Bir TensorRT motorunun ne \texttt{dict}'i, ne \texttt{.get()}'i,
        ne de veriye bağlı bir dallanması vardır.
  \item \textbf{SAM başlığı → bellek kodlayıcı} arasında bu karenin
        maskesi aynı \texttt{dict}'e geri yazılır.
  \item \textbf{Bellek dikkati → SAM başlığı} arasında ise hiçbir şey
        yoktur; bu sınır zorunlu değildir.
\end{itemize}

Değişken uzunluklu bellek sorunu, sabit 7 yuvaya doldurma (padding) ve
toplamsal bir dikkat maskesiyle işgal durumunu taşıma yoluyla
çözülmüştür -- bunun değişken-uzunluklu yolla birebir eşdeğer olduğu
testle kanıtlanmıştır (\texttt{test\_padded\_memory\_matches\_variable\_}
\texttt{length}). Motora taşınamayan şey, \emph{hangi} karelerin bu
yuvaları dolduracağına karar vermektir -- çünkü bu karar bir sözlük
sorgusu ve bir sıralamadır.

\subsection{Tek motora birleştirme mümkün müydü?}

İlke olarak evet: birleşik bir motor \texttt{(görüntü, doldurulmuş banka,
banka konumu, banka maskesi)} girdisini alıp \texttt{(maskeler, nesne
göstericisi, yeni bellek, yeni bellek konumu)} döndürebilirdi. Banka
derleme ve geri yazma yine motorun dışında kalırdı, çünkü bunlar bir
Python \texttt{dict}'i okuyup yazar.

Birleştirmenin kazandırabileceği tek şey, motor çağrıları arasındaki
\textbf{başlatma boşluğudur} (launch gap) -- bir motorun dönüşü ile
bir sonrakinin başlatılması arasındaki süre. Bu karar varsayıma değil
doğrudan ölçüme bırakılmıştır; sonucu ve eşik kuralı
\cref{sec:fuzyon-karari}'de sunulmaktadır.

\section{Sabit Şekilli Bellek Bankası ve CUDA Graph}

Bir CUDA Graph, kaydedildiği kernel çağrılarını \emph{aynı} cihaz
adreslerine karşı tekrar oynatır -- bu nedenle her motor kalıcı girdi/çıktı
arabellekleri tutar, adreslerini TensorRT'e bir kez verir, ısınır (TensorRT
bazı bağlam kurulumlarını ilk \texttt{enqueue}'a erteler), yakalar ve o
andan sonra her kare yalnızca arabelleklere kopyalama ile
\texttt{graph.replay()} çağrısından ibarettir.

Bunun çalışabilmesi için bellek bankasının \textbf{sabit şekilli}
olması gerekir. Gerçek kontrol noktası üzerinde izlenen anahtar
uzunluğu (\texttt{tools/analyze\_cost.py --trace-frames 24}), ilk 23
karede \textbf{16 farklı} değer almıştır (1. karede 516 belirteç,
7. karede 3612, 16. kareden sonra sabit 3648) -- her biri kendi CUDA
Graph yakalamasını, kendi TensorRT bağlamını gerektirirdi. Çözüm,
7 sabit yuvalı bir arabellek artı toplamsal bir dikkat maskesidir: boş
yuvalar $-30000$ ile maskelenir, ki bu değer hem fp16 hem fp32'de
$\exp(-30000) = 0$ altına taşarak softmax'ı canlı önek üzerinde
hesaplananla bit-eşdeğer kılar (gerçek $-\infty$ yerine sonlu bir değer
seçilmesinin nedeni, bir satırın tamamen maskelenmesi durumunda NaN
üretmemektir).

Bu doldurmanın zararsız olmasını sağlayan iki özellik: RoPE'nin
yuva-konumundan bağımsız olması (aynı frekans tablosu her yuvaya
kopyalanır, hangi yuvaların dolu olduğu tabloyu değiştirmez) ve zamansal
konum kodlamasının çağıran tarafından zaten \texttt{memory\_pos} içine
gömülmüş olması (yalnızca kopyalanır, yeniden hesaplanmaz).

\section{fp16 Hassasiyet Seçimi}
\label{sec:hassasiyet-secimi}

Motorlar fp16'da derlenmiştir, bf16'da değil -- ve bu bir uzlaşma değil,
Orin'in donanımının doğrudan sonucudur: bf16 aynı tensör çekirdeklerini
aynı hızda kullanır ama fp16'nın 11 bit anlamlı basamağına karşı yalnızca
8 bit taşır, ve TensorRT'in bf16 için daha az taktiği vardır. bf16'nın tek
avantajı olan geniş dinamik aralık, yalnızca etkinleşimler (activations)
fp16'yı taşırsa değer kazanır -- ölçüldü, taşmıyorlar
(\cref{sec:hassasiyet}). INT8'in neden \emph{motor olarak} kapsam dışı
bırakıldığı \cref{sec:int8-gelecek}'te ele alınmaktadır.

\section{Ölçüm Metodolojisi}
\label{sec:olcum-yontemi}

Bu projede birden fazla ölçüm aracı vardır ve her biri farklı bir soruya
cevap verir; ikisini karıştırmak yanıltıcı sonuçlara yol açar. Bu bölüm
o ayrımı sabitler.

\subsection{Bir karenin süresi neyi kapsar}
\label{sec:kare-tanimi}

Gerçek zamanlı bütçeye giren üç aşama, girmeyen bir aşama vardır:

\begin{table}[H]
  \centering
  \caption{Bir karenin zamanının bileşenleri.}
  \label{tab:kare-bilesenleri}
  \begin{tabular}{lll}
    \toprule
    Aşama & Ne & Bütçede mi? \\
    \midrule
    \texttt{pre} & kırpma + model girdisine ölçekleme + normalize etme & Evet \\
    \texttt{inference} & dört modül + EdgeTAM'ın Python defter tutması & Evet \\
    \texttt{post} & maskeleri kaynak çözünürlüğe geri ölçekleme & Evet \\
    \texttt{read} & kare dosyasını diskten okuma/çözme & \textbf{Hayır} -- ölçülür, hariç tutulur \\
    overlay + mp4 & maskeyi çizip video kodlama & \textbf{Hayır} -- yalnızca demo \\
    \bottomrule
  \end{tabular}
\end{table}

Bu tanım iki gerçek hatanın düzeltilmesiyle ortaya çıkmıştır. Birincisi,
\texttt{pre} başlangıçta ön-işleme değil, kaynak karenin görselleştirme
için diskten \emph{tekrar} okunmasıydı -- gerçek ön-işleme
\texttt{init\_state} içinde toplu yapılıyor, kurulum aşamasında kayboluyor
ve hiçbir kare süresinde görünmüyordu. İkincisi, video kodlama başlangıçta
ölçülen bölgenin içindeydi; artık akış halinde yazılıp ayrı ölçülmektedir.

\subsection{\texttt{read}'in bütçeden çıkarılması: gerekçe ve koşul}
\label{sec:read-disarida}

Ön-işleme ilk kez doğru ölçüldüğünde 1024 çözünürlükte modelin kendisi
kadar (\textasciitilde30~ms) pahalı çıkmıştır -- ayrıntılı çöküm
\cref{sec:onisleme-analizi}'dedir. Bunun ezici çoğunluğu (\textasciitilde35~ms)
görüntü dosyasının \emph{çözülmesi}dir, diskten okunması değil (aynı
dosyada ayrıştırıldığında: bayt okuma 0,5~ms, çözme 34,4~ms). Bu ayrım
önemlidir, çünkü ikisi farklı koşullarda ortadan kalkar: hiçbir dağıtım
diskten TIFF okumaz, ama \emph{sıkıştırılmış} kare alan bir dağıtım yine
de bir çözme işlemi öder.

Bu nedenle \texttt{read} artık overlay/kodlama ile aynı muameleyi görür:
ölçülür, raporlanır, bütçeye yazılmaz. Ancak bu çıkarımın geçerliliği
\textbf{koşulludur} ve rapor bunu açıkça belirtmelidir: kareler bellekte
ham geldiğinde yerine geçen maliyet ihmal edilebilir düzeydedir (bir
arabelleği sarmalamak $\sim$0,001~ms, bir kez kopyalamak $\sim$1,4~ms);
sıkıştırılmış geldiğinde ise bir çözme maliyeti geri döner -- daha ucuz,
ama sıfır değil. Bu projenin hedef dağıtımı (Ethernet üzerinden ham kare)
henüz kurulmadığından, bu koşul \cref{sec:ethernet-gelecek}'te açık bir
gelecek çalışma olarak bırakılmıştır.

\subsection{İki ``kare süresi'', ikisi de gerçek}

\begin{table}[H]
  \centering
  \caption{İki farklı ölçüm aracının cevapladığı farklı sorular.}
  \label{tab:iki-kare-suresi}
  \begin{tabular}{p{3.2cm}p{5.3cm}p{4.8cm}}
    \toprule
    Araç & Neyi ölçer & Ne için kullanılır \\
    \midrule
    \texttt{benchmark\_tracking.py} & yalnızca model; kareler baştan
      toplu çözülmüş & motor kümelerini veya çözünürlükleri birbirine
      karşı karşılaştırmak \\
    \texttt{cli.py} / \texttt{run\_records.py} & ön-işleme + model +
      son-işleme, kare başına & \textbf{dağıtılabilir kare bütçesi} \\
    \bottomrule
  \end{tabular}
\end{table}

Aradaki fark kare başına ön-işlemedir. Birini etiketsiz alıntılamak,
aynı sistemin 37,88~FPS ve 23,0~FPS gibi çelişkili görünen -- aslında
farklı şeyleri ölçen -- iki sayı vermesine yol açmıştır. Bu raporda her
sayı hangi araçtan geldiği belirtilerek sunulur.

\subsection{Yeniden Üretilebilirlik Tehlikeleri}
\label{sec:reprod-tehlike}

\begin{itemize}
  \item \textbf{Sabitlenmemiş saat hızları $\pm$\%30 fark yaratır.} Aynı
        yapılandırma iki ayrı koşuda 26,40~ms ve 33,40~ms ölçülmüştür.
        Rapora giren her sayıdan önce
        \texttt{sudo nvpmodel -m 0 \&\& sudo jetson\_clocks} çalıştırılmalı,
        ve \texttt{tools/capture\_environment.py} bu durumu kaydetmelidir.
  \item \textbf{İki adım aynı klibi kullanmalı.} Aksi halde biri diğerinden
        farklı bir içerik üzerinde ölçüm yapar ve karşılaştırma anlamsızlaşır.
  \item \textbf{\texttt{--offload-video} tüm araçlarda tutarlı olmalı.}
        Aksi halde EdgeTAM her kareyi GPU'da fp32 olarak tutar --
        500 kare için \textasciitilde6,3~GB -- ve Orin'in paylaşımlı
        belleğinde bu modelin kendisini yavaşlatır.
  \item \textbf{Biriktirilmiş kareler modelle bellek bant genişliği için
        yarışır.} CPU ve GPU, Orin'de belleği \emph{ve} bant genişliğini
        paylaşır; işlenmiş klibi bir listede tutmak, modelin bağımlı
        olduğu kaynağı tüketmektir.
\end{itemize}
